\documentclass[10pt]{article}

% ---------- Document Identity (update these only) ----------
\newcommand{\DocStage}{}
\newcommand{\DocTitle}{Your Road to Tier One SOF}
\newcommand{\ProgramName}{182-Week Civilian-to-Selection Readiness Pipeline}
\newcommand{\ProgramSubtitle}{}

% ---------- Page ----------
\usepackage[legalpaper,left=0.5in,right=0.5in,top=0.6in,bottom=0.45in]{geometry}

% ---------- Typography ----------
\usepackage[T1]{fontenc}
\usepackage[utf8]{inputenc}
\usepackage{libertinus}
\usepackage{microtype}
\usepackage{parskip}
\usepackage{ragged2e}
\usepackage{textcomp}
\AtBeginDocument{\color{black}}
\AtBeginDocument{\pagecolor{white}}
\AtBeginDocument{\justifying}

% ---------- Landscape pages ----------
\usepackage{pdflscape}

% ---------- Color ----------
\usepackage[table]{xcolor}
\definecolor{StageBlue}{HTML}{1F4E79}
\definecolor{StageBlueDark}{HTML}{163A5A}
\definecolor{LightGray}{HTML}{F2F4F7}
\definecolor{MidGray}{HTML}{D9DEE5}
\definecolor{RuleGray}{HTML}{C7CED8}
\definecolor{HeaderInk}{HTML}{000000}
\definecolor{Ink}{HTML}{000000}

% ---------- Utility ----------
\usepackage{enumitem}
\sloppy
\setlength{\emergencystretch}{2em}

% ---------- Boxes / UI ----------
\usepackage[most]{tcolorbox}

\tcbset{
  charterTitle/.style={
    enhanced,
    colback=StageBlue!4,
    colframe=StageBlue,
    boxrule=1.25pt,
    arc=3pt,
    left=16pt,right=16pt,top=14pt,bottom=14pt,
    borderline north={3.2pt}{0pt}{StageBlue},
    borderline west={4pt}{0pt}{StageBlue},
    borderline south={1.25pt}{0pt}{StageBlue},
    drop shadow={black!25!white},
  },
  charterRule/.style={
    enhanced,
    colback=LightGray,
    colframe=RuleGray,
    boxrule=0.85pt,
    arc=3pt,
    left=12pt,right=12pt,top=10pt,bottom=10pt,
    borderline west={3pt}{0pt}{StageBlue},
  },
  charterBadge/.style={
    enhanced,
    colback=StageBlue,
    coltext=white,
    colframe=StageBlueDark,
    boxrule=0.6pt,
    arc=2pt,
    left=6pt,right=6pt,top=3pt,bottom=3pt,
  }
}
\newtcbox{\Badge}{nobeforeafter,charterBadge}

% ---------- Lists ----------
\setlist[itemize]{leftmargin=1.5em, itemsep=0.25em, topsep=0.25em}
\setlist[enumerate]{leftmargin=1.8em, itemsep=0.25em, topsep=0.25em}

% ---------- TOC ----------
\usepackage{tocloft}
\setcounter{tocdepth}{3}
\setlength{\cftbeforesecskip}{3pt}
\setlength{\cftbeforesubsecskip}{2pt}
\setlength{\cftbeforesubsubsecskip}{0pt}
\renewcommand{\cftsecfont}{\bfseries}
\renewcommand{\cftsecpagefont}{\bfseries}
\renewcommand{\cftsubsecfont}{\normalfont}
\renewcommand{\cftsubsecpagefont}{\normalfont}
\renewcommand{\cftsubsubsecfont}{\normalfont}
\renewcommand{\cftsubsubsecpagefont}{\normalfont}
\setlength{\cftsecindent}{0pt}
\setlength{\cftsubsecindent}{1.25em}
\setlength{\cftsubsubsecindent}{2.8em}
\setlength{\cftsecnumwidth}{2.1em}
\setlength{\cftsubsecnumwidth}{2.8em}
\setlength{\cftsubsubsecnumwidth}{3.6em}

% Remove the built-in TOC heading so only the custom heading is shown
\makeatletter
\renewcommand{\tableofcontents}{\@starttoc{toc}}
\makeatother

% ---------- Header/Footer: page number only ----------
\usepackage{fancyhdr}
\pagestyle{fancy}
\fancyhf{}
\renewcommand{\headrulewidth}{0pt}
\renewcommand{\footrulewidth}{0pt}
\fancyfoot[C]{\thepage}

% ---------- Section formatting ----------
\usepackage{titlesec}

\titleformat{\section}
  {\Large\bfseries\color{Ink}}
  {\thesection}{0.6em}{}
  [\vspace{0.25em}\textcolor{StageBlue}{\rule{\linewidth}{0.8pt}}]

\titleformat{\subsection}
  {\large\bfseries\color{Ink}}
  {\thesubsection}{0.6em}{}

\titleformat{\subsubsection}
  {\normalsize\bfseries\color{Ink}}
  {\thesubsubsection}{0.6em}{}

\titlespacing*{\section}{0pt}{0.95em}{0.35em}
\titlespacing*{\subsection}{0pt}{0.65em}{0.25em}
\titlespacing*{\subsubsection}{0pt}{0.55em}{0.2em}

% ---------- Table engine ----------
\usepackage{tabularray}
\UseTblrLibrary{booktabs}

% ---------- tabularray: GLOBAL table treatment ----------
\SetTblrInner{
  cells = {fg=black, bg=white, valign=t},
}

% ---------- Header helpers ----------
\newcommand{\Hdr}[1]{\shortstack[c]{\strut #1\strut}}

% ---------- Lines ----------
\newcommand{\TitleLine}{\textcolor{StageBlue}{\rule{0.98\linewidth}{0.95pt}}}
\newcommand{\SubTitleLine}{\textcolor{RuleGray}{\rule{0.98\linewidth}{0.65pt}}}

% ---------- Continued on next page ----------
\newcommand{\ContinuedOnNextPageInline}{%
  \par
  \vfill
  \noindent\textcolor{RuleGray}{\rule{\linewidth}{1pt}}\vspace{-2pt}
  \noindent\hfill\textit{Continued on next page}\par\vspace{-10pt}
  \noindent\textcolor{RuleGray}{\rule{\linewidth}{1pt}}
  \clearpage
}

% ---------- PDF hyperlinks ----------
\usepackage[hidelinks]{hyperref}
\hypersetup{
  colorlinks=false,
  pdfborder={0 0 0},
  linktoc=all,
  pdftitle={\DocTitle},
  pdfauthor={\ProgramName}
}

% =========================================================
\begin{document}

\section*{Stage 1.F — Governance Rules and Validity Doctrine}
\phantomsection
\addcontentsline{toc}{section}{Stage 1.F — Governance Rules and Validity Doctrine}

\subsection*{1.F.1 Governance Hierarchy}
\phantomsection
\addcontentsline{toc}{subsection}{1.F.1 Governance Hierarchy}

Ordered hierarchy (highest authority first):

\begin{enumerate}
\item Stage 0 overrides downstream content. If any downstream rule conflicts with Stage 0 constraints/boundaries, the downstream rule is non-authoritative until corrected and re-versioned.
\item Safety boundaries override performance goals and schedule. When safety triggers occur, \textbf{STOP} \textbf{NOW} actions apply immediately; evidence from unsafe continuation is inadmissible.
\item Governance overrides programming preferences. If a planned session cannot meet required elements/admissibility rules, it must be tagged accurately (\textbf{MODIFIED}/\textbf{INVALID}) rather than forcing compliance by unsafe means.
\item Validity rules override “good intentions.” If evidence cannot be verified (missing minimum dataset, missing required elements, or boundary violations), it is inadmissible for gate decisions.
\end{enumerate}

Evidence produced under compromised safety (stop-rule continuation, water-governance breach, falsified/incomplete documentation) \textbf{MUST} be treated as inadmissible for advancement and \textbf{MUST} trigger default conservative gate posture (\textbf{HOLD} or \textbf{MEDICAL} \textbf{HOLD} as applicable).

\newpage

\subsection*{1.F.2 Session Validity Doctrine}
\phantomsection
\addcontentsline{toc}{subsection}{1.F.2 Session Validity Doctrine}

\subsubsection*{1.F.2.1 Required Session Elements (Non-Negotiable)}
\phantomsection

Required elements checklist (order is locked):

\begin{itemize}
\item Safety self-check
\item Warm-up
\item Mobility
\item Main work
\item Cardio
\item Cooldown
\item Recovery
\item Post-session notes
\end{itemize}

Body text enforcement (not embedded in checklist labels):

\begin{itemize}
\item Cardio occurs in every training session.
\item Cardio minimum standard is minimum 30 minutes continuous per training session.
\item If Cardio is under the minimum or is not continuous, the session is \textbf{INVALID}  (no gate credit).
\end{itemize}

Minimum logging fields by element (governance-level; no programming prescriptions):

Safety self-check (\textbf{MUST} log):

\begin{itemize}
\item readiness snapshot (sleep hours since last session; general fatigue level),
\item pain flags (location + 0–10 severity),
\item red-flag symptom check (presence/absence; details only if present),
\item environment check (surface hazard, temperature band, wind/ice flag),
\item water exposure check (if any aquatic facility is involved: certified supervision present Y/N; underwater/breath-hold planned Y/N).
\end{itemize}

Warm-up (\textbf{MUST} log):

\begin{itemize}
\item completed Y/N,
\item modality summary (e.g., dynamic movement, easy walk),
\item duration (minutes).
\end{itemize}

Mobility (\textbf{MUST} log):

\begin{itemize}
\item completed Y/N,
\item target regions (shoulder/hip/ankle/general),
\item duration (minutes),
\item pain or restriction notes if relevant.
\end{itemize}

Main work (\textbf{MUST} log):

\begin{itemize}
\item completed Y/N,
\item domain(s) performed (strength / calisthenics / run / ruck / swim / durability),
\item modality summary (what was done at a high level),
\item duration (minutes),
\item intensity anchor (RPE and/or talk-test cue if applicable).
\end{itemize}

Cardio (\textbf{MUST} log):

\begin{itemize}
\item modality,
\item minutes (must be ≥ the minimum stated above),
\item intensity anchor: RPE (1–10) + talk-test cue,
\item continuity statement: “continuous completed” or “not continuous” (if not continuous → \textbf{INVALID}).
\end{itemize}

Cooldown (\textbf{MUST} log):

\begin{itemize}
\item completed Y/N,
\item modality summary (easy walk, stretch),
\item duration (minutes).
\end{itemize}

Recovery (\textbf{MUST} log):

\begin{itemize}
\item sleep hours,
\item soreness 0–10,
\item fatigue 0–10,
\item pain flags,
\item foot hotspots,
\item recovery actions taken at category level only.
\end{itemize}

Post-session notes (\textbf{MUST} log):

\begin{itemize}
\item brief factual summary,
\item what changed vs plan,
\item substitutions,
\item constraints for next session,
\item \textbf{SESSION-RC} reason code if \textbf{MODIFIED} or \textbf{INVALID}.
\end{itemize}

\subsubsection*{1.F.2.2 Tag Definitions (Operational)}
\phantomsection

\textbf{VALID}

\begin{itemize}
\item All required elements completed.
\item Cardio completed and meets the minimum: minimum 30 minutes continuous per training session.
\item Minimum logging dataset complete for the session (element completion + core fields).
\item No stop-rule continuation occurred (\textbf{STOP} \textbf{NOW} triggers were not ignored).
\end{itemize}

\textbf{MODIFIED}

\begin{itemize}
\item All required elements completed.
\item Cardio completed and meets the minimum continuous standard.
\item A documented substitution/downgrade occurred to preserve intent while reducing risk.
\item The substitution \textbf{MUST} NOT increase injury risk beyond current tolerance.
\item Minimum logging dataset complete AND includes:
\begin{itemize}
\item trigger/signal prompting modification,
\item what changed (planned vs executed),
\item \textbf{SESSION-RC} reason code.
\end{itemize}
\end{itemize}

\textbf{INVALID} 
Any of the following forces \textbf{INVALID} :

\begin{itemize}
\item Any required element missing (not completed).
\item Cardio under the minimum standard or not continuous.
\item Required logging dataset missing or insufficient to verify compliance.
\item Undocumented deviation such that intent or compliance cannot be verified.
\item \textbf{STOP} \textbf{NOW} trigger occurred and the session continued in a manner that compromised safety governance (evidence becomes inadmissible).
\end{itemize}

Admissibility for gates:

\begin{itemize}
\item \textbf{VALID} sessions are admissible.
\item \textbf{MODIFIED} sessions are admissible only within the cap rules defined in 1.F.4 and \textbf{MUST} use \textbf{SESSION-RC}.
\item \textbf{INVALID} sessions are not admissible and \textbf{MUST} NOT be used to justify progression.
\end{itemize}

\subsubsection*{1.F.2.3 Reason Codes (Required; \textbf{SESSION-RC})}
\phantomsection

\textbf{SESSION-RC} codes \textbf{MUST} be used for any \textbf{MODIFIED} or \textbf{INVALID} session. Codes are assigned to the primary driver of the deviation (one primary code minimum; a secondary code MAY be added if needed for audit clarity).

\textbf{SESSION-RC} Code Set (Canonical; Expandable Via Patch Without Weakening Doctrine):

\begin{itemize}
\item \textbf{SESSION-RC-001} (\textbf{ENV}) — Applies to: \textbf{MODIFIED} / \textbf{INVALID} — Environment hazard prevents safe execution (ice, extreme wind, unsafe surface, air quality).
\begin{itemize}
\item Minimum documentation: hazard type; location; date/time; substitute used; intensity anchor; safety self-check notes.
\item Default gate posture if repeated: if ≥2 in window → \textbf{HOLD}; require robust indoor substitute plan.
\end{itemize}
\item \textbf{SESSION-RC-002} (PAIN) — Applies to: \textbf{MODIFIED} / \textbf{INVALID} — Pain or discomfort alters mechanics or indicates unsafe loading tolerance.
\begin{itemize}
\item Minimum documentation: pain location + 0–10; mechanics change observed Y/N; substitution details; continuity status.
\item Default gate posture if repeated: \textbf{HOLD} or \textbf{REGRESS}; consider \textbf{MEDICAL} \textbf{HOLD} if stop-rule thresholds triggered.
\end{itemize}
\item \textbf{SESSION-RC-003} (\textbf{ILL}) — Applies to: \textbf{MODIFIED} / \textbf{INVALID} — Illness symptoms reduce safe training tolerance.
\begin{itemize}
\item Minimum documentation: symptoms category; RPE increase trend; substitution or termination; safety self-check.
\item Default gate posture if repeated: default \textbf{HOLD}; \textbf{MEDICAL} \textbf{HOLD} if stop-rule categories triggered.
\end{itemize}
\item \textbf{SESSION-RC-004} (\textbf{EQUIP}) — Applies to: \textbf{MODIFIED} / \textbf{INVALID} — Equipment access failure prevents planned modality (facility closed; gear unavailable).
\begin{itemize}
\item Minimum documentation: what failed; substitute modality; duration; intensity anchor; continuity statement.
\item Default gate posture if repeated: \textbf{HOLD} until the dependency is addressed and documented.
\end{itemize}
\item \textbf{SESSION-RC-005} (\textbf{TIME}) — Applies to: \textbf{MODIFIED} / \textbf{INVALID} — Time constraint prevents completion of required elements (including Cardio minimum).
\begin{itemize}
\item Minimum documentation: which elements shortened/missed; why; minutes completed; whether Cardio minimum met.
\item Default gate posture if repeated: if Cardio minimum missed → \textbf{INVALID}; repeated \textbf{TIME} failures → \textbf{HOLD} with corrective actions.
\end{itemize}
\item \textbf{SESSION-RC-006} (\textbf{TRAVEL}) — Applies to: \textbf{MODIFIED} / \textbf{INVALID} — Travel disrupts environment/equipment access.
\begin{itemize}
\item Minimum documentation: location; substitute modality; minutes; intensity anchor; daily steps status.
\item Default gate posture if repeated: default \textbf{HOLD} if repeated; must preserve admissible sessions via substitution ladder.
\end{itemize}
\item \textbf{SESSION-RC-007} (\textbf{SAFETY}) — Applies to: \textbf{MODIFIED} / \textbf{INVALID} — Safety self-check triggers conservative downgrade or termination (non-medical governance trigger).
\begin{itemize}
\item Minimum documentation: trigger/signal; immediate action; documentation; notifications.
\item Default gate posture if repeated: default \textbf{HOLD}; \textbf{MEDICAL} \textbf{HOLD} if stop-rule categories triggered.
\end{itemize}
\item \textbf{SESSION-RC-008} (\textbf{WEATHER}) — Applies to: \textbf{MODIFIED} / \textbf{INVALID} — Weather shifts mid-session require immediate change (storm, extreme cold/heat onset).
\begin{itemize}
\item Minimum documentation: weather shift; action taken; substitute; environment fields.
\item Default gate posture if repeated: \textbf{HOLD}; enforce pre-session hazard planning.
\end{itemize}
\item \textbf{SESSION-RC-009} (\textbf{ADMIN}) — Applies to: \textbf{INVALID} — Logging failure or missing required documentation (cannot verify compliance).
\begin{itemize}
\item Minimum documentation: what was missing; whether verifiable reconstruction from same-day records was attempted; why it failed.
\item Default gate posture if repeated: repeated \textbf{ADMIN} invalidity → \textbf{HOLD} until logging compliance restored.
\end{itemize}
\item \textbf{SESSION-RC-010} (\textbf{WATER}) — Applies to: \textbf{INVALID} — Any water governance breach (unauthorized underwater or breath-hold attempt; lack of certified supervision).
\begin{itemize}
\item Minimum documentation: what occurred; supervision status; immediate stop; notifications; corrective actions.
\item Default gate posture if repeated: immediate \textbf{MEDICAL} \textbf{HOLD} posture; governance breach requires revalidation and corrective control.
\end{itemize}
\end{itemize}

Application rule:

\begin{itemize}
\item \textbf{SESSION-RC} is not a justification to “count” non-compliant work. It is an audit tool to describe reality and enforce conservative gate posture when trends repeat.
\end{itemize}

\newpage

\subsection*{1.F.3 Test Validity Doctrine}
\phantomsection
\addcontentsline{toc}{subsection}{1.F.3 Test Validity Doctrine}

Tests are tagged \textbf{VALID} or \textbf{INVALID} only.

Governance Requirements for a Test to Be Tagged \textbf{VALID} (Protocol-Agnostic):
A test \textbf{MUST} be tagged \textbf{VALID} when all of the following are true:

\begin{itemize}
\item The test references the adopted test identity (\textbf{ID}/name) used by the phase row (protocol specifics exist elsewhere).
\item Conditions are standardized enough to compare results:
\begin{itemize}
\item route or course consistency (or clearly logged deviation),
\item equipment consistency (footwear, ruck load configuration, pool lane length where applicable),
\item timing and measurement tool adequacy (timer/GPS/measured course status logged).
\end{itemize}
\item Safety posture is preserved:
\begin{itemize}
\item no stop-rule triggers were ignored,
\item \textbf{STOP} \textbf{NOW} triggers end the attempt.
\end{itemize}
\item Warm-up was completed and logged (completed Y/N + duration).
\item Environment/log context is recorded at a high level:
\begin{itemize}
\item course measured Y/N (or pool lane length),
\item surface category,
\item weather category,
\item wind/ice flag,
\item temperature band,
\item supervision status for any aquatic setting (certified supervision present Y/N).
\end{itemize}
\end{itemize}

A test \textbf{MUST} be tagged \textbf{INVALID} when any of the following occur:

\begin{itemize}
\item Core documentation missing such that conditions cannot be verified.
\item Measurement integrity is compromised (unknown course distance; unknown load; unknown lane length).
\item Stop-rule event occurred and the attempt continued such that safety governance was compromised.
\item Water governance is violated (any underwater or breath-hold attempt without certified supervision).
\end{itemize}

Corrective action posture:

\begin{itemize}
\item If a test is gate-required and is \textbf{INVALID}, default gate posture is \textbf{HOLD} until the test is rerun under \textbf{VALID} conditions (unless other disqualifiers require \textbf{REGRESS} or \textbf{MEDICAL} \textbf{HOLD}).
\end{itemize}

Test Validity Table:

\begin{itemize}
\item \textbf{VALID} — Meets documentation, safety, and condition standardization requirements; timing/measurement and environment recorded — Gate impact: admissible evidence.
\item \textbf{INVALID} — Missing documentation, compromised measurement integrity, stop-rule governance breach, or water governance breach — Corrective action: rerun under controlled conditions; correct measurement/logging; review stop-rule adherence — Gate impact: not admissible; default \textbf{HOLD} when gate-required.
\end{itemize}

\textbf{TEST-RC} (Reason Codes For \textbf{INVALID} Tests; Separate From \textbf{SESSION-RC}):
\textbf{TEST-RC} codes \textbf{MUST} be used whenever a test is tagged \textbf{INVALID}. One primary code minimum.

\begin{itemize}
\item \textbf{TEST-RC-001} (\textbf{MEAS}) — Measurement integrity failure (course not measured; load not confirmed; lane length unknown) — Default gate posture: \textbf{HOLD} until valid retest.
\item \textbf{TEST-RC-002} (\textbf{ENV}) — Environment not recorded or conditions materially inconsistent without documentation — Default gate posture: \textbf{HOLD} until valid retest.
\item \textbf{TEST-RC-003} (\textbf{PROTO}) — Protocol reference missing or execution non-comparable to adopted standard — Default gate posture: \textbf{HOLD} until valid retest.
\item \textbf{TEST-RC-004} (\textbf{SAFETY}) — \textbf{STOP} \textbf{NOW} trigger occurred or safety governance was compromised — Default gate posture: \textbf{HOLD} or \textbf{MEDICAL} \textbf{HOLD} depending on trigger category.
\item \textbf{TEST-RC-005} (\textbf{ADMIN})  — Documentation failure (time not recorded; missing test sheet; lost record) — Default gate posture: \textbf{HOLD}; repeated \textbf{ADMIN} failures may force \textbf{REGRESS}.
\item \textbf{TEST-RC-006} (\textbf{WATER}) — Water governance violation — Default gate posture: \textbf{MEDICAL} \textbf{HOLD}; governance-breach handling required.
\end{itemize}

Evidence artifacts note:

\begin{itemize}
\item Video, GPS, and photos are recommended (not mandatory).
\item Missing recommended artifacts \textbf{MUST} NOT automatically invalidate a test unless core validity requirements above are violated.
\end{itemize}

\newpage

\subsection*{1.F.4 Gate Decision Doctrine}
\phantomsection
\addcontentsline{toc}{subsection}{1.F.4 Gate Decision Doctrine}

\subsubsection*{1.F.4.1 Definitions}
\phantomsection

\begin{itemize}
\item \textbf{ADVANCE} — Authorization to progress based on admissible evidence meeting thresholds and absence of disqualifiers.
\item \textbf{HOLD} — Maintain current phase/progression level due to insufficient admissible evidence or unresolved risk posture.
\item \textbf{REGRESS} — Step back to safer exposure level when tolerance is exceeded or evidence indicates instability.
\item \textbf{MEDICAL} \textbf{HOLD} — Progression paused/restricted due to stop-rule categories or clinician-imposed limitations (governance posture only).
\end{itemize}

\subsubsection*{1.F.4.2 Evidence Requirements (Default Rules)}
\phantomsection

\begin{itemize}
\item Default decision window: two most recent complete training weeks preceding the gate week (must not be shorter than two complete weeks absent patch).
\item \textbf{ADVANCE} eligibility: ≥5/6 admissible sessions per week on average across the decision window.
\item \textbf{MODIFIED} cap: outside medical issues, ≤1 \textbf{MODIFIED} per week admissible.
\item Binary disqualifier: Any \textbf{INVALID} session within the decision window disqualifies \textbf{ADVANCE}.
\item Gate-required tests must be \textbf{VALID}; an \textbf{INVALID} gate-required test defaults to \textbf{HOLD} until a valid retest.
\item Dependencies for next phase must be validated; risk flags/trends must be reviewed.
\end{itemize}

\subsubsection*{1.F.4.3 Gate Table (Required)}
\phantomsection

\begin{itemize}
\item \textbf{ADVANCE} — requires a ≥5/6 admissible average; \textbf{MODIFIED} cap compliant; zero \textbf{INVALID}; gate-required tests \textbf{VALID}; dependencies validated — disqualifiers include any \textbf{INVALID}, any required test \textbf{INVALID}, unresolved stop-rule events, unresolved gait-altering mechanics, water governance violation.
\item \textbf{HOLD} — default when advancement cannot be justified; corrective stabilization and retest posture as required.
\item \textbf{REGRESS} — trend indicates exposure exceeds tolerance; step back and re-govern progression.
\item \textbf{MEDICAL} \textbf{HOLD} — stop-rule category requires clinician-led clearance or restrictions or risk posture is unacceptable.
\end{itemize}

\subsection*{1.F.5 Stop Rules and Escalation Posture (Aligned to Stage 0.5)}
\phantomsection
\addcontentsline{toc}{subsection}{1.F.5 Stop Rules and Escalation Posture (Aligned to Stage 0.5)}

Stop rules apply to training sessions and test events. The default action is \textbf{STOP} \textbf{NOW}. Where indicated by category, default gate posture becomes \textbf{MEDICAL} \textbf{HOLD}.

Minimum \textbf{STOP} \textbf{NOW} documentation: trigger, time/context, action taken, escalation path, outcome and return-to-training condition, resulting validity tag.

\subsection*{1.F.6 Freeze Windows (Change Limits During Execution)}
\phantomsection
\addcontentsline{toc}{subsection}{1.F.6 Freeze Windows (Change Limits During Execution)}

\subsubsection*{1.F.6.1 Definition and Default Within-Phase Model}
\phantomsection

Freeze window: execution-affecting changes are prohibited to preserve comparability and prevent drift.

Default model:

\begin{itemize}
\item Week 1: ramp-in adjustment permitted (within Stage 0 constraints); document.
\item Weeks 2 -> mid-phase test week: frozen (only safety substitutions).
\item Post mid-phase test -> end-phase test week: frozen (only safety substitutions).
\item Final weeks: frozen (no structural changes; only safety substitutions).
\end{itemize}

\subsubsection*{1.F.6.2 Allowed vs Patch-Required Changes}
\phantomsection

Safety substitutions and controlled day swaps may be allowed; changes affecting intent, gate tests, thresholds, cadence, required elements, phase duration, or high-consequence exposures are patch-required and must not be executed silently.

\subsection*{1.F.7 Change Control (Patch Protocol + Change Log Schema)}
\phantomsection
\addcontentsline{toc}{subsection}{1.F.7 Change Control (Patch Protocol + Change Log Schema)}

\subsubsection*{1.F.7.1 Patch Rules}
\phantomsection

Patch must include: what, why, where, effective date, validation check, owner, impact level, revalidation scope.
Stage 0 changes are constraint changes and trigger full revalidation including Stage 8 crosswalk rerun before downstream artifacts remain deployable.

\subsubsection*{1.F.7.2 Patch Record Template}
\phantomsection

\begingroup
\scriptsize
\setlength{\tabcolsep}{2pt}
\renewcommand{\arraystretch}{1.12}

\begin{longtblr}{
  width=\linewidth,
  colspec = {
    X[l,.5]
    X[l,.5]
    X[l,.5]
    X[l,.5]
    X[l,.5]
    X[l,.5]
    X[l,.5]
    X[l,.65]
  },
  hlines,
  vlines,
  cells = {hsep=6pt, vsep=6pt, halign=l, valign=t},
  row{odd} = {bg=white},
  row{even} = {bg=LightGray},
  row{1} = {bg=StageBlue!15, fg=black, font=\bfseries, halign=l, valign=m},
}
\Hdr{Patch \textbf{ID}} & \Hdr{Date} & \Hdr{Target} & \Hdr{Change} & \Hdr{Rationale} & \Hdr{Validation} & \Hdr{Owner} & \Hdr{Impact Level} \\
\end{longtblr}

\endgroup

\subsection*{1.F.8 Minimum Dataset for Logging (Decision-Grade Evidence)}
\phantomsection
\addcontentsline{toc}{subsection}{1.F.8 Minimum Dataset for Logging (Decision-Grade Evidence)}

Per-session, per-day, and per-week logs must be sufficient for a third-party reviewer to determine admissibility, threshold attainment, safety compliance, and whether gate decisions were justified.
\newpage
\subsection*{1.F.9 Deviation Handling (How to Correct Drift Without Hiding It)}
\phantomsection
\addcontentsline{toc}{subsection}{1.F.9 Deviation Handling (How to Correct Drift Without Hiding It)}

\begin{itemize}
\item Mis-tags must be corrected with original tag, corrected tag, reason, correction date, and Coach Lead acknowledgment.
\item \textbf{INVALID} gate-required tests must be rerun under controlled conditions; retest posture prioritizes standardization and safety over peak-output chasing.
\item Environmental changes must be recorded; if comparability is not preserved, treat results cautiously and require additional validation exposures.
\item Missed weeks default \textbf{HOLD}; no make-up volume; conservative re-entry; \textbf{MEDICAL} \textbf{HOLD} when stop-rule categories apply.
\end{itemize}

\subsection*{1.F.10 Conflict-Resolution Rule}
\phantomsection
\addcontentsline{toc}{subsection}{1.F.10 Conflict-Resolution Rule}

\begin{itemize}
\item Any downstream artifact violating 1.F governance rules \textbf{MUST} be treated as Draft — Non-Deployable until corrected and re-versioned.
\item Gate decisions made using inadmissible evidence are invalid and \textbf{MUST} be rerun using admissible evidence within the decision-window rules.
\end{itemize}

\end{document}